\phantomsection
\chapter*{Empat Sahabat, Satu Keluarga}
\addcontentsline{toc}{section}{Empat Sahabat, Satu Keluarga --- Aliif Abdu Mukhti A.}
\penulis{Aliif Abdu Mukhti A.}


\awalcerpen{A}{da empat orang sahabat} yang sudah berteman sejak lama. Mereka adalah Aliif, Firman, Hilbam, dan Ikiy. Walaupun sekarang mereka bersekolah di tempat yang berbeda, persahabatan mereka tetap berjalan seperti biasa.

Aliif bersekolah di Pratama Adi. Firman bersekolah di Banjar Asri. Hilbam memilih untuk belajar di pesantren, sedangkan Ikiy bersekolah di Taruna.

Awalnya mereka sempat khawatir karena setelah masuk sekolah masing-masing, mereka tidak akan bisa sering bertemu seperti dulu.

``Kayaknya nanti kita bakal jarang ketemu,'' kata Firman.

``Jangan gara-gara beda sekolah terus jadi lupa teman,'' jawab Aliif.

Hilbam tersenyum dan berkata, ``Sekolah boleh beda, tapi pertemanan jangan ikut beda.''

Ikiy kemudian menimpali sambil bercanda, ``Pokoknya kalau ada yang lupa, nanti kita tagih!''

Mereka semua tertawa.

Sejak saat itu, mereka membuat kesepakatan untuk tetap menjaga persahabatan. Kalau sedang ada waktu luang, mereka akan menyempatkan diri untuk bertemu. Kalau tidak bisa bertemu, mereka tetap berkomunikasi melalui pesan.

Hari-hari pertama terasa sedikit berbeda. Aliif sibuk dengan kegiatan di Pratama Adi. Firman mulai memiliki teman-teman baru di Banjar Asri. Hilbam menjalani kehidupan yang berbeda di pesantren. Sementara Ikiy mulai terbiasa dengan kegiatan dan aturan di Taruna.

Walaupun begitu, mereka tidak pernah benar-benar melupakan satu sama lain.

Suatu malam, mereka sedang mengobrol melalui \textit{grup chat}.

``Gimana sekolah kalian?'' tanya Aliif.

``Lumayan, banyak tugas,'' jawab Firman.

``Kalau aku kegiatannya lumayan padat,'' kata Ikiy.

Hilbam kemudian membalas, ``Di pesantren juga banyak kegiatan. Tapi aku senang.''

Mereka saling bercerita tentang kehidupan di sekolah masing-masing. Dari obrolan sederhana itu, mereka sadar bahwa walaupun lingkungan mereka berbeda, mereka tetap memiliki satu kesamaan: mereka masih saling peduli.

Beberapa minggu kemudian, mereka akhirnya memiliki kesempatan untuk berkumpul.

Aliif datang lebih dulu. Tidak lama kemudian Firman datang, disusul Hilbam dan Ikiy.

``Wah, akhirnya lengkap lagi!'' teriak Firman.

Mereka langsung tertawa dan saling bercanda seperti tidak pernah lama berpisah.

Mereka kemudian duduk bersama dan bercerita tentang pengalaman masing-masing.

Aliif bercerita tentang kegiatan di Pratama Adi. Firman menceritakan teman-temannya di Banjar Asri. Hilbam bercerita tentang kehidupan di pesantren, sedangkan Ikiy menceritakan pengalaman selama sekolah di Taruna.

Mereka baru menyadari bahwa ternyata banyak hal yang berubah setelah mereka bersekolah di tempat yang berbeda.

Namun, ada satu hal yang tidak berubah.

Persahabatan mereka.

Suatu ketika, Firman sedang mengalami masalah dan merasa sangat tertekan dengan keadaan di sekolah. Ia tidak tahu harus bercerita kepada siapa.

Akhirnya ia menghubungi ketiga sahabatnya.

``Aku lagi banyak pikiran,'' kata Firman.

Aliif langsung menjawab, ``Cerita saja. Kita dengerin.''

Hilbam juga berkata, ``Jangan dipendam sendirian.''

Ikiy menambahkan, ``Kita memang beda sekolah, tapi kita tetap teman kamu.''

Firman merasa lega setelah menceritakan masalahnya. Ia sadar bahwa meskipun mereka tidak berada di sekolah yang sama, ia masih memiliki tiga orang yang selalu siap mendengarkan.

Waktu terus berjalan. Mereka semakin sibuk dengan sekolah masing-masing. Pertemuan yang dulunya bisa dilakukan hampir setiap hari sekarang hanya sesekali.

Namun, setiap kali bertemu, rasanya tetap sama.

Mereka masih bisa tertawa bersama, saling mengejek, dan saling membantu.

Suatu sore, mereka kembali berkumpul di tempat yang biasa mereka datangi sejak dulu.

``Lihat, sekarang kita sudah beda sekolah semua,'' kata Aliif.

``Iya, tapi masih lengkap'', jawab Hilbam.

Firman tersenyum.

``Berarti benar kata orang, saudara itu enggak harus satu sekolah.''

Ikiy kemudian berkata, ``Bahkan enggak harus satu keluarga juga.''

Mereka semua terdiam sebentar, lalu tersenyum.

Mereka memang tidak memiliki hubungan darah. Mereka juga tidak lagi berada di sekolah yang sama. Aliif di Pratama Adi, Firman di Banjar Asri, Hilbam di pesantren, dan Ikiy di Taruna.

Tetapi jarak dan perbedaan sekolah tidak membuat mereka berpisah.

Bagi mereka, persahabatan bukan tentang seberapa sering bertemu, tetapi tentang seberapa kuat mereka tetap saling peduli.

Empat teman yang berbeda sekolah akhirnya menjadi seperti empat saudara yang dipertemukan oleh persahabatan.

\jedahias
